% ============================================================
% ICS & AI Transparency Log --- Meichen Ding (2362344)
%
% This file is the single source for two deliverables:
%   (a) Appendix A of project2/meichending/Method.tex, where it is included
%       after \appendix with \def\reportbuild{} set beforehand, so that the
%       cross-references resolve to the numbers of that report;
%   (b) the standalone document project2/ics/MeichendingICS.tex, which inputs
%       this file without \reportbuild, so the same sentences appear with
%       words instead of numbers.
% Format follows Checklist/ICS_bvp.tex: Part A is the team AI transparency
% log, Part B is the individual contribution statement.
% ============================================================

% Cross-reference helper. In the report build the numbered objects exist and
% are printed; in the standalone ICS document they do not, so plain words are
% printed instead and the document has no unresolved references.
\ifdefined\reportbuild
  \newcommand{\reftext}[2]{#1}
\else
  \newcommand{\reftext}[2]{#2}
\fi

\section{AI Transparency Log}
\label{app:ai_log}

This appendix records the distance between raw generative-AI output and the
claims that survived into this report; it is a record of auditing, not a count
of prompts. Every claim listed in A2 was re-derived, re-run against the team's
own code, or falsified before it was allowed into the text, and each of the
three audits changed a claim that an earlier draft had already stated. The
conversation transcripts are retained by the team and can be supplied on
request; they are not appended here.

\subsection*{A1. AI tools used}

\begin{table}[H]
\centering\small
\caption{Generative-AI tools used in Project 2. Entries describe the use that
actually reached the report; routine autocomplete and spelling checks are
covered by the last row.}
\label{tab:ai_tools}
\begin{tabular}{@{}>{\raggedright\arraybackslash}p{3.9cm}>{\raggedright\arraybackslash}p{3.3cm}>{\raggedright\arraybackslash}p{7.4cm}@{}}
\toprule
Tool & Version / date & Primary use \\
\midrule
DeepSeek-based assistant in DeepSeek Harness (\texttt{deepseek-flash})
  & current build, Oct--Nov 2026
  & Drafting and re-deriving \S1.2 and \S2.1--\S2.4 (model, weak form,
    truncation error, interface conductance, boundary elimination, discrete
    structure); writing the independent verification scripts in
    \path{support/verify/}; adversarial review of the section text. \\
Claude (Anthropic), web application
  & current build, Oct--Nov 2026
  & Drafting of the algorithm section, integration notes for merging the
    section files, figure and table captions, wording and consistency passes. \\
ChatGPT (OpenAI) and Gemini (Google)
  & current builds, Oct--Nov 2026
  & \LaTeX{} error diagnosis, document structure, phrasing and completeness
    checks. \\
\midrule
All of the above, routine use
  & ---
  & Autocomplete, spelling, reference formatting, table alignment. No numerical
    value was accepted without reproduction (see A3). \\
\bottomrule
\end{tabular}
\end{table}

\subsection*{A2. Significant AI-assisted claim audits}

Three interactions are recorded below because each one changed a mathematical
or numerical claim rather than simply rewording it. The common failure pattern
is a claim that is true for the face coefficient, or for one dimension, being
silently promoted to the solution of the two-dimensional problem.

{\footnotesize
\begin{longtable}{@{}p{0.95cm}>{\raggedright\arraybackslash}p{7.4cm}>{\raggedright\arraybackslash}p{5.8cm}@{}}
\caption{Claim audits. Column 3 states what would have shown the claim to be wrong.}\label{tab:ai_audits}\\
\toprule
Item & Prompt / claim / assumptions & Risk and possible falsifier \\
\midrule
\endfirsthead
\caption[]{Claim audits (continued).}\\
\toprule
Item & Prompt / claim / assumptions & Risk and possible falsifier \\
\midrule
\endhead
\bottomrule
\endfoot

1 &
``Arithmetic averaging of the interface conductance produces an
$\mathcal{O}(1)$ error in the interface flux, so the arithmetic-averaged
solution is $\mathcal{O}(1)$ inaccurate.''
\textbf{Assumption hidden in the wording:} the face-coefficient error and the
error of the discrete solution are the same object. &
The claim mixes three different quantities: the relative face-conductance
error $(1-c)^2/(4c)$ (genuinely $\mathcal{O}(1)$ in $h$), the nodal error and
the midline interface-flux error (first order in $h$), and the \emph{maximum}
interface-flux error over the square, which is $\mathcal{O}(1)$ for a different
reason --- the exact Dirichlet trace on $y=0,1$. \textbf{Falsifier:} a
refinement study. If the nodal error did not halve under mesh refinement, the
first-order part of the claim would be wrong. \\
2 &
``Since $A_h \succeq a_{\min}A_h^{(1)}$ in the Loewner order and both matrices
are $M$-matrices, the inverses compare entrywise,
$A_h^{-1}\le a_{\min}^{-1}(A_h^{(1)})^{-1}$; hence
$\|A_h^{-1}\|_\infty\le 1/(8a_{\min})$.'' \textbf{Assumption hidden:}
spectral order implies entrywise order. &
The implication is false: monotonicity of $M$-matrix inverses needs an
\emph{entrywise} hypothesis, while growing $a$ makes the off-diagonals
$-a_f/h^2$ \emph{more negative}. \textbf{Falsifier:} evaluate both inverses for
a strongly varying coefficient, e.g.\ $a=1+9\sin^2(\pi x)$ at $n=63$, where the
entrywise ratio reaches $3.10>1$. The conclusion
$\|A_h^{-1}\|_\infty\le1/(8a_{\min})$ is nevertheless true and is therefore
proved by a maximum-principle barrier argument instead. \\
3 &
``The arithmetic operator is symmetric and its shared face flux is
conservatively balanced, so the interface physics is correct.'' &
Symmetry and local conservation are necessary but not sufficient; they are
satisfied by a matrix whose interface flux is wrong by $36\%$.
\textbf{Falsifier:} compare the computed interface flux with the exact
$J_*=-20/11$ on a benchmark whose exact solution is known. It differs by
$0.662$ while the two-sided reconstructed jump is $3.85\times10^{-13}$ and the
symmetry defect is zero. \\
\end{longtable}}


\subsection*{A3. Independent verification and decision}

{\footnotesize
\begin{longtable}{@{}p{0.95cm}>{\raggedright\arraybackslash}p{6.2cm}>{\raggedright\arraybackslash}p{7.0cm}@{}}
\caption{What was checked, where the evidence lives, and what changed.}\label{tab:ai_verification}\\
\toprule
Item & Independent check and evidence location & Final decision / modification / limitation \\
\midrule
\endfirsthead
\caption[]{Independent verification (continued).}\\
\toprule
Item & Independent check and evidence location & Final decision / modification / limitation \\
\midrule
\endhead
\bottomrule
\endfoot

1 &
Re-derived the truncated-series argument in one dimension, where the discrete
system reduces to a tridiagonal chain with a closed-form flux; reproduced the
leading term $h(a_L-a_R)^2/(a_L+a_R)^2$ analytically and confirmed it against
$N=31,\dots,255$ (\path{support/verify/check_interface_order.py},
\reftext{Table~\ref{tab:interface_refinement}}{\S4.3}). Then re-assembled the
full two-dimensional layered problem and compared the nodal, midline and
maximal face-flux errors with \S4.3. &
The two-layer statement of \S2.2 and
\reftext{Remark~\ref{rem:arith_2d_boundary}}{the remark on the two-dimensional
boundary} replaced the single blanket $\mathcal{O}(1)$ claim: first-order nodal
and midline flux error, $\mathcal{O}(1)$ maximal face-flux error near the
prescribed horizontal boundaries, $\mathcal{O}(1)$ face-conductance error.
\textbf{Limitation:} the maximum is a boundary effect of the Dirichlet square;
it is stated as an observation on this benchmark, not as an estimate valid for
other boundary data. \\
2 &
Ran the counterexample search and the bound check in
\path{support/verify/check_stability.py}: entrywise ratio for
$a=1+9\sin^2(\pi x)$, $\|A_h^{-1}\|_\infty$ against $1/(8a_{\min})$ for five
coefficients in one and two dimensions, and the comparison with the
constant-coefficient maximum $0.0736<\frac18$
(\reftext{Remark~\ref{rem:entrywise_false}}{the remark recording the
counterexample}, \reftext{Remark~\ref{rem:conv_constant}}{the remark on the
constant}). &
The proof of \reftext{Theorem~\ref{thm:convergence}}{the convergence theorem}
was rewritten with a maximum-principle barrier: $A_h\bm\varphi=\bm 1$ exactly
for $\varphi=\frac12x(1-x)/a_{\min}$ extended by its own boundary values, and
$A_h(\bm\varphi-\bm w)=0$ with $\bm\varphi-\bm w\ge0$ on $\partial\Omega$.
\textbf{Limitation:} equality in $\|\bm w\|_\infty\le1/(8a_{\min})$ holds only
in one dimension; in two the constant is an upper bound ($0.0736$ versus
$0.125$). \\
3 &
Separated the three diagnostics on the same computed solution: symmetry defect
$\|A_h-A_h^T\|/\|A_h\|$, reconstructed two-sided flux jump, and error against
the exact flux
(\reftext{\S\ref{sec:results}}{\S4} and the interface diagnostics of
\reftext{\S\ref{sec:algorithm_impl}}{\S3}; independent re-computation in
\path{test/validate_interface.py}). &
\reftext{Remark~\ref{rem:symmetry_vs_accuracy}}{The remark on symmetry versus
accuracy} and \S4.3 now report the three numbers side by side instead of
treating symmetry as evidence of physical accuracy. \textbf{Limitation:} the
jump diagnostic is only defined where a two-sided reconstruction exists, so it
is not a substitute for an exact solution. \\
\end{longtable}}


Two further routine checks are recorded for completeness. The smooth benchmark
is \emph{blind} to the $\theta$-correction of the face coefficient when
$a=1+x$: arithmetic and midpoint averaging then produce entrywise identical
matrices, so the second-order table in \S4.2 cannot validate the correction
term (\reftext{Remark~\ref{rem:theta_blind}}{the remark on the
$\theta$-correction}). And no convergence claim is made anywhere in this report
without a named norm and a four-level refinement table; claims about the
layered benchmark are exactness statements
(\reftext{Theorem~\ref{thm:nodal_exactness}}{the nodal-exactness theorem}), not
order statements, because $u_*\notin H^2(\Omega)$ there
(\reftext{Remark~\ref{rem:thm2_scope}}{the remark on the validity of the
truncation analysis}).

\subsection*{A4. Intellectual ownership declaration}

The mathematical formulation, the discretisation, the solvers, the experiments,
the validation evidence and the conclusions are the team's own work.
Generative AI was used for drafting, clarification and preliminary checking, and
every AI-assisted claim was audited as recorded above. Each member affirms that
they can explain and modify the team-owned discretisation and solver code,
justify the mathematical claims in their own sections, and state the
assumptions and the tests applied at each library interface, without AI
assistance during the viva. The viva may ask for a small live modification or
diagnostic on team-owned work.

\begin{table}[H]
\centering\small
\caption{Ownership declaration. This copy carries the signature of Meichen
Ding; the other four rows are completed and signed by each member in their own
copy of this table. A declaration is never drafted on another member's
behalf.}
\label{tab:ai_ownership}
\begin{tabular}{@{}>{\raggedright\arraybackslash}p{2.7cm}p{1.7cm}>{\raggedright\arraybackslash}p{4.9cm}>{\raggedright\arraybackslash}p{4.6cm}@{}}
\toprule
Name & Signature & Team-owned artefact / contribution & Trusted dependency and interface \\
\midrule
Meichen Ding (2362344)
  & \underline{\hspace{1.2cm}}
  & \S1.2; \S2.1--\S2.4; the remarks on the $\theta$-correction, on mesh
    parity, on the two-dimensional boundary and on symmetry versus accuracy;
    Appendix~A; the scripts in \path{support/verify/}
  & \texttt{numpy.linalg} dense solves and
    \texttt{scipy.sparse.linalg.eigsh} used only for corroboration: checked
    against explicit residuals, against the exact layered solution and against
    a dense Cholesky factorisation \\
\bottomrule
\end{tabular}
\end{table}

\section{Individual Contribution Statement}
\label{app:ics}

Part A above is the team log. The statement below is individual: it states in
the first person what I did, and it cross-references the audits of
Section~\ref{app:ai_log} where the corresponding AI-assisted claim was checked.

\begin{formbox}\small
\textbf{Individual Contribution Statement --- Project 2 (PDE BVP)}

\textbf{Name:} Meichen Ding (2362344)

\textbf{Role this project:} Mathematical theory lead (consistency, adaptivity
and convergence): owner of the model formulation, the discrete scheme, its
truncation-error and interface analysis, and the discrete-structure results.

\textbf{The specific things I did this project:}
\begin{enumerate}[leftmargin=1.4em,itemsep=0.15em]
  \item
  Wrote \S1.2 and \S2.1--\S2.4 of the report's \path{Method.tex}: the continuous
  problem and its hypotheses, the flux-balance derivation that fixes
  one shared face coefficient per edge, the index map and block-tridiagonal
  structure of the assembled operator (with the non-zero count $5n^2-4n$), the
  symmetric Dirichlet elimination and the hand-computed $4\times4$ matrix at
  $n=2$ used as the small-case oracle, the layered exact solution, the
  second-order local truncation error with its $\theta$-correction, the
  series-resistance proposition for the interface conductance with the relative
  error $(1-c)^2/(4c)$, the SPD / $M$-matrix / discrete-maximum-principle
  theorem, and the discrete stability and $\mathcal{O}(h^2)$ convergence theorem
  with its maximum-principle barrier argument.
  \item
  Wrote the independent verification scripts \path{support/verify/}:
  \texttt{check\_interface\_order.py} (one-dimensional conservative scheme,
  $a_L=1$, $a_R=10$, $N=31,63,127,255$; harmonic nodes exact to $10^{-14}$ and
  every face flux equal to $-20/11$, arithmetic nodal error first order with
  measured orders $1.015/1.008/1.004$ and the closed-form chain flux
  reproduced), \texttt{check\_stability.py} (the entrywise-ratio counterexample
  at $a=1+9\sin^2(\pi x)$, ratio $3.10$; $\|A_h^{-1}\|_\infty$ against
  $1/(8a_{\min})$ for five coefficients) and
  \texttt{check\_matrix\_evidence.py} (symmetry defect, \texttt{eigsh} smallest
  eigenvalue and dense Cholesky on the submitted assembler). These produced the
  numbers in the appendix tables and in \S2.4.
  \item
  Ran the AI-assisted claim audits of Part A: the three-layer separation of the
  arithmetic-averaging errors (audit 1), the falsification of the
  Loewner-implies-entrywise argument and its replacement by the barrier proof
  (audit 2), and the separation of ``symmetric and conservative'' from
  ``physically correct'' (audit 3), including the numerical values that decide
  each one ($0.662$ versus $3.85\times10^{-13}$; $3.10>1$; $0.0736$ versus
  $0.125$).
\end{enumerate}

\textbf{What I would do differently:}
I would pin down the mesh convention with the code owner before writing any
convergence sentence. The skeleton uses $h=1/(n+1)$ with $n$ interior points
while the report uses $h=1/N$ with $N$ intervals, so the same symbol $N$ means
different things in the skeleton, in the layered tables and in the FEM
comparison; the interface-parity requirement $N$ odd for the FDM and $N$ even
for the FEM then has to be stated as design rather than discovered late. I would
also have demanded the two-dimensional layered numbers before writing the
one-dimensional flux estimate, because that is exactly where the two-layer
statement became necessary.

\textbf{Anything I should have asked for help with earlier:}
I should have asked, in Week 5, for each member to commit their own sections and
their own ICS from their own account. By the time the appendices were written
the repository history was a single upload commit, so authorship had to be
reconstructed from memory and from section ownership instead of being visible in
\texttt{git log}. I also should have asked for a review of the stability proof
before it went into a draft: the Loewner-order step survived an internal
re-reading and was only caught when it was written out as a counterexample
search.

\vspace{0.8em}
\textbf{Signature:} \underline{\hspace{5cm}} \hfill
\textbf{Date:} \underline{\hspace{3cm}}
\end{formbox}
